% Fichier généré automatiquement par tools/generate_corrections.py.
% Source : GEO/GEO-03/12_BielleManivelle/12_BielleManivelle.tex
% Statut : complet (5/5 question(s) renseignée(s), 0 reprise(s)).
\begin{corrigebox}[Corrigé extrait de la banque]
\CorrigeQuestion{1}
\begin{marginfigure}
\includegraphics[width=\linewidth]{12_01_c}
\end{marginfigure}
\CorrigeQuestion{2}
En réalisant une fermeture géométrique, on obtient 
$\vect{AB}+\vect{BC}+\vect{CA}=\vect{0}$ 
$ \Leftrightarrow R\vect{i_1} - L\vect{i_2} - \lambda(t) \vect{j_0} = \vect{0}$.
On projette alors cette expression dans $\rep{0}$ : 

$
\left\{
\begin{array}{l}
R\cos\theta(t)- L\cos\varphi(t) = 0 \\
R\sin\theta(t) - L\sin\varphi(t) - \lambda(t) = 0 
\end{array}
\right.
$.

On cherche à éliminer $\varphi(t)$ : 

$
\left\{
\begin{array}{l}
R\cos\theta(t) = L\cos\varphi(t)  \\
R\sin\theta(t)  - \lambda(t) = L\sin\varphi(t)
\end{array}
\right.
$.

En élevant au carré, on a donc 

$\left\{
\begin{array}{l}
R^2\cos^2\theta(t) = L^2\cos^2\varphi(t)  \\
\left(R\sin\theta(t)  - \lambda(t)\right)^2 = L^2\sin^2\varphi(t)
\end{array}
\right.
$.

En conséquence, 
$R^2\cos^2\theta(t)  + \left(R\sin\theta(t)  - \lambda(t)\right)^2 = L^2$ et 

$ \left(R\sin\theta(t)  - \lambda(t)\right)^2 = L^2 - R^2\cos^2\theta(t) $ 
$\Rightarrow    \lambda(t) = \pm\sqrt{L^2 - R^2\cos^2\theta(t)} + R\sin\theta(t) $ .
\CorrigeQuestion{3}
$\lambdap(t) = \pm\left( \dfrac{R^2 \thetap(t) \cos\theta(t)\sin\theta(t)}{\sqrt{L^2 - R^2\cos^2\theta(t)}}\right) + \thetap(t)R\cos\theta(t) $
\CorrigeQuestion{4}
On retient la branche positive de la position :
\[
\lambda(t)=\sqrt{L^2-R^2\cos^2\!\theta(t)}+R\sin\theta(t).
\]
Avec $\theta(t)=\omega t$ et $\omega=\SI{100}{rad.s^{-1}}$, on obtient
\[
\dot\lambda(t)=\omega\left(
\frac{R^2\cos(\omega t)\sin(\omega t)}
{\sqrt{L^2-R^2\cos^2(\omega t)}}
+R\cos(\omega t)
\right).
\]
La courbe est calculée pour $R=\SI{10}{mm}$, puis pour
$L=\SI{20}{mm}$ et $L=\SI{30}{mm}$.

\begin{marginfigure}
\includegraphics[width=\linewidth]{12_02_c}
\end{marginfigure}
\CorrigeQuestion{5}
L'accélération peut être calculée par différence centrée à partir de la vitesse :
\[
\ddot\lambda(t_i)\simeq
\frac{\dot\lambda(t_{i+1})-\dot\lambda(t_{i-1})}{2\,\Delta t}.
\]
\begin{marginfigure}
\includegraphics[width=\linewidth]{12_03_c}
\end{marginfigure}

\begin{enumerate}
\item Tracer le graphe des liaisons.
\item $\lambda(t) = \pm\sqrt{L^2 - R^2\cos^2\theta(t)} + R\sin\theta(t)$.
\item $\lambdap(t) = \pm\left( \dfrac{R^2 \thetap(t) \cos\theta(t)\sin\theta(t)}{\sqrt{L^2 - R^2\cos^2\theta(t)}}\right) + \thetap(t)R\cos\theta(t)$.
\item Évaluer la vitesse pour les deux valeurs de $L$.
\item Dériver numériquement la vitesse pour obtenir l'accélération.
\end{enumerate}
\end{corrigebox}
